\chapter*{Preface: A Player's Companion}

\addcontentsline{toc}{chapter}{Preface}

\textit{``The coin that never spends is the one you don't remember taking.''} \\
−−− Serafine of the Velvet Touch

Welcome, player. This book is not a spell to be memorised, nor a contract to be obeyed. It is a companion – something to keep at your elbow when the dice are rolling and the story is racing ahead. You will not need every page tonight. What you will need is a clear sense of how your character works, a quick way to look up that one rule you forgot, and permission to put the book down and play.

\subsection*{How This Book Is Organised}

\begin{itemize}
    ℹtem \textbf{\autoref{ch:intro}} – The philosophy of narrative‑first play, safety tools, and Session Zero. Read this before you build a character.
    ℹtem \textbf{\autoref{ch:core−mechanics}} – The 90‑second loop, Position, Difficulty, Boons, Story Beats, Harm, and Fatigue. This is the engine. Keep the Outcome Matrix (\autoref{sec:outcome−matrix}) and the Position table (\autoref{sec:position}) close at hand.
    ℹtem \textbf{\autoref{ch:character−creation}} – Building your first character (30−−36XP). Use the checklist in \autoref{sec:bonds−complications} to get started in ten minutes.
    ℹtem \textbf{\autoref{ch:advancement}} – Attributes, Skills, dice pools, Experience Points, and the Tier system (Rookie to Mythic). You will return to this chapter every time you earn XP. Bookmark \autoref{sec:spending−xp} and the training times in \autoref{sec:standard−training}.
    ℹtem \textbf{\autoref{ch:talents}} – Minor, Major, Prestige, and Epic talents. Skim the minor talents first; the rest will grow with your character.
    ℹtem \textbf{\autoref{ch:magic}} – All five paths (Free Caster, Runekeeper, Invoker, Cantor, Summoner). If you are new, read only \autoref{sec:free−caster} (Free Caster) or \autoref{sec:runekeeper} (Runekeeper). Save the others for a second character.
    ℹtem \textbf{\autoref{ch:healing}} – How to recover Fatigue and Harm, with and without magic. Keep the ℎyperref[sec:healing−summary]{Healing at a Glance} table bookmarked.
    ℹtem \textbf{\autoref{ch:assets−followers}} – Safehouses, spy networks, loyal companions, and terrestrial patrons. You will not need this until Tier II (41+ XP), but when you do, \autoref{sec:asset−free−use} and \autoref{sec:using−followers} are the sections to remember.
    ℹtem \textbf{\autoref{ch:currency−free−economy}} – No coin counting. Use XP, Strings, and Debt Timers instead. \autoref{sec:no−coin} gives you the mindset; the rest is reference.
    ℹtem \textbf{\autoref{ch:world−cultures}} – The peoples and regions of Fate's Edge. Read only the culture your character comes from. The rest is for when you travel.
    ℹtem \textbf{\autoref{ch:world−powers−obligation}} – The Patrons (Carrion‑King, Traveler, Oath of Flame, etc.) and their Rites, Obligation, and Corruption. This is a deep dive. For a first character, stick to the six recommended Patrons in \autoref{subsec:choosing−a−patron}.
    ℹtem \textbf{\autoref{ch:example−concepts}} – Ten example characters (Guardian, Mage, Rogue, etc.). Use them as inspiration or as pre‑gens.
    ℹtem \textbf{\autoref{ch:enhanced−play}} – Tactical advice, momentum banking, optional rules like war mounts and character death. Come back here when you want to push your play.
    ℹtem \textbf{\autoref{ch:downtime}} – What you do between adventures: recovery, training, crafting, carousing. The Rush Rule (\autoref{sec:rush−rule}) is a favourite.
    ℹtem \textbf{\autoref{ch:legacy−engine}} – How your character's story can continue through a successor. Read this only when you are thinking about retirement or a heroic death.
\end{itemize}

\subsection*{A Reading Plan for New Players}

You do not need to read 200 pages to play your first session. Here is what you actually need:

\begin{enumerate}
    ℹtem Read \autoref{ch:intro} – especially the safety tools (Lines, Veils, X‑Card) and the Session Zero checklist.
    ℹtem Read \autoref{ch:core−mechanics} up to \autoref{sec:boons} – the play loop, Position, Difficulty, and the Outcome Matrix.
    ℹtem Skim \autoref{sec:harm−fatigue} to understand how Harm and Fatigue work. You will refer back to it.
    ℹtem Use the character creation checklist (\autoref{ch:character−creation}) to build your first character. The pre‑generated examples in \autoref{ch:example−concepts} can speed this up.
    ℹtem If you are playing a magical character, read \autoref{ch:first−magical}. It walks you through a Runekeeper step by step.
    ℹtem Print the one‑page player reference (if your GM provides it) or copy the Outcome Matrix onto an index card.
    ℹtem Close the book and play.
\end{enumerate}

\subsection*{What to Bookmark for Your Second Session}

After you have rolled dice once, you will want these sections close at hand:

\begin{itemize}
    ℹtem \textbf{Harm and Fatigue (\autoref{sec:take−harm−fatigue})} – The procedure for armour conversion and the ``roller‑coaster effect'' (taking Harm clears Fatigue) is easy to forget.
    ℹtem \textbf{Spending Boons (\autoref{sec:spending−boons})} – You can re‑roll a die, improve Position, or activate an Asset.
    ℹtem \textbf{HEAL Cantrip (\autoref{subsec:HEAL−cantrip})} – The fastest way to clear Fatigue.
    ℹtem \textbf{Skill synergies (\autoref{sec:skill−synergies})} – When you can use a different Attribute with your skill, the game opens up.
    ℹtem \textbf{The Rush Rule (\autoref{sec:rush−rule})} – When you want to train faster and risk a flaw.
\end{itemize}

\subsection*{A Note on the Quickstart}

If you are completely new, start with the \textit{Fate's Edge Quickstart}. It contains the core rules, four pre‑generated characters, and a complete adventure (``The Lantern at Dusk''). Run that first. Then open this book to build your own character for the \textit{Blood and Silk} trilogy.

\subsection*{The Golden Rule (for Players)}

\begin{quote}
``When in doubt, make the choice that serves the story. If you find yourself asking `What's the optimal move?' try reframing it as `What would be most interesting for the story?'''
\end{quote}

That permission is written into the rules. Use it.

\subsection*{Acknowledgements}

This book exists because players like you asked ``What if?'' and GMs like yours answered ``Let's find out.'' Thank you for sitting at the table, for sharing your character's fears and hopes, and for making the story worth telling.

Now build a character who owes a debt, carries a scar, and has a name that someone will remember. The road is waiting.

\vspace{1em}
\noindent\textit{−−− The Grey Wanderer} \\
\textit{Castra Ferrum, Year of the Iron Harvest}